\section{Estructura de la tesis}

En el Capítulo \ref{marco_teorico} se aporta una explicación detallada de todos los conceptos abordados en el marco de este trabajo así como la explicación de los conceptos fundacionales que dan lugar a todas las técnicas. Así, se parte de qué son las imágenes, cómo se forman, qué operaciones se pueden realizar sobre ellas y cómo es que dan lugar a tareas como la visión por computadora y SLAM. Además, se presentan los grupos y álgebras de Lie, así como su integración en SLAM y en el modelo de Gaussian Splatting. Del mismo modo, se describe toda la matemática que resulta necesaria para comprender los sensores y el modelo de Gaussian Splatting trabajados, así como las técnicas de optimización utilizadas.

En el Capítulo \ref{estado_del_arte} se presenta un análisis de los trabajos más relevantes en los que se basa este trabajo.

En el Capítulo \ref{metodo} se brinda una breve explicación de la arquitectura del programa, así como las decisiones de diseño tomadas, complementando la explicación teórica con la forma computacional del algoritmo. Además, se describen de forma minimal los conceptos básicos necesarios para comprender la implementación.

En el Capítulo \ref{resultados} se muestran los resultados obtenidos comparando este trabajo con otros algoritmos de SLAM existentes en la literatura, utilizando distintos conjuntos de datos y las métricas comunes en SLAM.

En el Capítulo \ref{conclusiones} se brindan las conclusiones obtenidas a partir de los resultados, así como las limitaciones del trabajo y las posibles líneas de investigación futura.
